\documentclass[11pt,a4paper]{article}

% Encoding and fonts
\usepackage[utf8]{inputenc}
\usepackage[T1]{fontenc}
\usepackage{lmodern}
\usepackage{microtype}
\exhyphenpenalty=50
\tolerance=2000
\emergencystretch=1.5em
\sloppy
\setlength{\parskip}{0.6em}
\setlength{\parindent}{1.5em}

% Page geometry
\usepackage[margin=1.8cm]{geometry}

% URL line breaking
\usepackage[hyphens]{url}

% Links
\usepackage{hyperref}
\hypersetup{
  colorlinks=true,
  linkcolor=blue!60!black,
  citecolor=blue!60!black,
  urlcolor=blue!60!black
}

% Pandoc Syntax Highlighting Polyfills
\usepackage{framed}
\definecolor{shadecolor}{RGB}{248,248,248}
\newenvironment{Shaded}{\begin{snugshade}}{\end{snugshade}}
\newenvironment{Highlighting}[]{\ttfamily\raggedright}{\par}
\newcommand{\NormalTok}[1]{#1}
\newcommand{\KeywordTok}[1]{\textbf{#1}}
\newcommand{\DataTypeTok}[1]{#1}
\newcommand{\DecValTok}[1]{#1}
\newcommand{\BaseNTok}[1]{#1}
\newcommand{\FloatTok}[1]{#1}
\newcommand{\ConstantTok}[1]{#1}
\newcommand{\CharTok}[1]{#1}
\newcommand{\SpecialCharTok}[1]{#1}
\newcommand{\StringTok}[1]{#1}
\newcommand{\VerbatimStringTok}[1]{#1}
\newcommand{\SpecialStringTok}[1]{#1}
\newcommand{\ImportTok}[1]{#1}
\newcommand{\CommentTok}[1]{\textit{#1}}
\newcommand{\DocumentationTok}[1]{\textit{#1}}
\newcommand{\AnnotationTok}[1]{\textit{#1}}
\newcommand{\CommentVarTok}[1]{\textit{#1}}
\newcommand{\OtherTok}[1]{#1}
\newcommand{\FunctionTok}[1]{#1}
\newcommand{\VariableTok}[1]{#1}
\newcommand{\ControlFlowTok}[1]{\textbf{#1}}
\newcommand{\OperatorTok}[1]{#1}
\newcommand{\BuiltInTok}[1]{#1}
\newcommand{\ExtensionTok}[1]{#1}
\newcommand{\PreprocessorTok}[1]{#1}
\newcommand{\AttributeTok}[1]{#1}
\newcommand{\RegionMarkerTok}[1]{#1}
\newcommand{\InformationTok}[1]{\textit{#1}}
\newcommand{\WarningTok}[1]{\textit{#1}}
\newcommand{\AlertTok}[1]{\textbf{\color{red}#1}}
\newcommand{\ErrorTok}[1]{\textbf{\color{red}#1}}

% Math
\usepackage{amsmath,amssymb}

% Tables
\usepackage{booktabs}
\usepackage{longtable}
\usepackage{array}
\usepackage{calc}  % CRITICAL: defines \real{} for pandoc columns
\usepackage{float}
\renewcommand{\arraystretch}{1.3}

% Graphics
\usepackage{graphicx}
\usepackage{xcolor}

% Pandoc compatibility: define \tightlist
\providecommand{\tightlist}{%
  \setlength{\itemsep}{0pt}\setlength{\parskip}{0pt}}

\title{\textbf{Intent of Thought (IoT): A Governance Lifecycle for LLM Reasoning Topology Selection}}

\author{
  Naveen Riaz Mohamed Kani \\
  Independent Researcher, Melbourne, Australia \\
  ORCID: \href{https://orcid.org/0009-0003-9173-2425}{0009-0003-9173-2425} \\
  \texttt{naveenriaz@synthai.biz}
}

\date{March 2026}

\begin{document}
\maketitle

\begin{abstract}
Large language models can reason through chain-of-thought, tree-of-thought, and graph-of-thought topologies, yet no established mechanism governs which topology to deploy for a given task or how to recover when the wrong one is selected. We introduce Intent of Thought (IoT), a pre-reasoning governance layer that specifies task intent through a triple of Purpose, Anti-Purpose, and Success Signal, then uses this specification to guide topology selection, monitor reasoning alignment, and diagnose failure. 

\par\medskip

We formalise the IoT Lifecycle: a closed loop from intent capture (five modes of varying fidelity) through topology selection, drift detection, retrospective judgement (classifying failures as False Capture, False Selection, or False Execution), and learning. We evaluate IoT governance across 705 scored reasoning episodes, 7 models (7B to 120B parameters), 18 tasks, and 4 topologies. 

\par\medskip

Adding an IoT governance triple improves reasoning quality by +15.7\% (n = 486). Weaker models benefit disproportionately (+33\% for 7B vs. +16\% for 120B), suggesting governance functions as a cognitive equaliser. Governance rescues mismatched topologies: tree-of-thought quality improves +37\% under IoT. At the model scales tested, chain-of-thought is a robust default across all task categories; the framework's empirical value is in governance uplift and failure protection rather than optimal topology selection. A longitudinal case study of a multi-agent debate protocol that independently evolved IoT-like governance over 75 iterations validates the lifecycle in a production setting.

\par\medskip\noindent\textbf{Keywords:} Large Language Models, Chain-of-Thought Reasoning, Reasoning Topology, Intent Governance, Failure Modes, Multi-Agent Debate
\end{abstract}

\newpage
\tableofcontents
\newpage

\section{Introduction}\label{introduction}

The proliferation of reasoning topologies for large language
models, from simple chain-of-thought prompting {[}Wei et al., 2022{]}
through tree-of-thought {[}Yao et al., 2023{]} and graph-of-thought
{[}Besta et al., 2024{]}, has unlocked unprecedented capabilities, yet
simultaneously created a critical governance gap. Each topology imposes
a strict structural commitment on how a model reasons: CoT constrains
reasoning to a linear derivation chain, ToT introduces branching and
backtracking, and GoT permits structural merging and refinement across
nodes. However, as AI systems transition from passive, short-horizon
question answering to active, multi-step autonomous agency, there
remains no established mechanism to govern \emph{which} topology should
be deployed for a given task, \emph{what} specific failure vectors the
reasoning must avoid, or \emph{how} to structurally recover when an
agent confidently executes the wrong cognitive pattern.

This gap has severe practical consequences. When reasoning topology is
decoupled from task intent, models frequently experience ``catastrophic
drift'', where the structural execution is flawless, but the semantic
trajectory diverges entirely from the user's actual goal. In our
experimental evaluation across 486 scored reasoning episodes (Section
5), using tree-of-thought on interconnected analytical tasks produced a
mean quality score of 1.00 out of 3.00, a catastrophic failure. The
same tasks under chain-of-thought scored 2.47. The topology itself was
not defective; rather, the \emph{governance} connecting the task's true
intent to its structural execution was absent. Practitioners currently
face a silent hazard: relying on default topologies that are
structurally sound but strategically wrong. The core problem for the
next generation of agentic systems is not merely predicting a single
`optimal' topology, but providing the structural governance necessary to
prevent these catastrophic reasoning mismatches before they compound.

We introduce \textbf{Intent of Thought (IoT)},\footnote{We use IoT
  throughout to denote Intent of Thought, not to be confused with the
  Internet of Things or Iteration of Thought {[}Radha et al., 2024{]}.}
a pre-reasoning governance layer deeply inspired by
Belief-Desire-Intention (BDI) agent architectures but distinctly adapted
for the stochastic nature of LLMs. IoT addresses this governance gap
through three primary contributions:

\begin{enumerate}
\def\labelenumi{\arabic{enumi}.}
\item
  \textbf{The IoT Framework} (Section 3). We define an explicit
  \emph{intent triple} comprising a stated Purpose (the desired
  outcome), an Anti-Purpose (the specific boundary conditions the
  reasoning must avoid), and a Success Signal (the verification criteria
  used to judge adequacy). A topology selection function maps this
  triple to a recommended cognitive structure, while a continuous drift
  detection score monitors alignment during active inference.
\item
  \textbf{The IoT Lifecycle} (Section 4). Traditional frameworks assume
  intent is given correctly and governance operates forward-only. We
  relax these assumptions by introducing the \emph{Capture Spectrum}
  (five escalating modes of intent elicitation, ranging from zero-shot
  inference to collaborative, multi-turn learned capture) and
  \emph{Retrospective Judgement} (a backward intent reconstruction
  protocol that diagnoses systemic failure as False Capture, False
  Selection, or False Execution). Together, these mechanisms form a
  continuously compounding closed-loop lifecycle: Capture \textgreater{}
  Select \textgreater{} Monitor \textgreater{} Judge \textgreater{}
  Learn \textgreater{} Capture.
\item
  \textbf{Empirical Evidence} (Section 5). We rigorously evaluate IoT
  governance across 705 scored reasoning episodes spanning 7 models
  (from 7B open weights to 120B proprietary parameters), 18 high-stakes
  tasks across 6 domains, and 4 reasoning topologies. Key empirical
  findings include:

  \begin{itemize}
  \tightlist
  \item
    \textbf{Baseline Governance Uplift}: Structuring reasoning through
    an IoT governance triple improves aggregate reasoning quality by
    \textbf{+15.7\%} (p \textless{} 0.001, N = 486).
  \item
    \textbf{The Governance Equaliser}: Weaker models benefit
    disproportionately from explicit governance. 7B parameter models
    gained up to +33.3\% in reasoning quality, while 120B models gained
    +16.3\%, demonstrating that intent scaffolding acts as a critical
    cognitive equaliser for open-weight deployments.
  \item
    \textbf{The Capability Ceiling Effect}: Governance uplift exhibits
    diminishing returns for state-of-the-art frontier models (e.g.,
    Claude 3.5 Haiku, GPT-4o-mini). This indicates that while advanced
    proprietary models increasingly internalise intent-to-topology
    mappings, smaller models fundamentally require the explicit
    guardrails IoT provides.
  \item
    \textbf{Protection Against Mismatch}: Governance successfully
    rescues mismatched topologies. Default tree-of-thought quality
    improves from 1.67 to 2.28 (+37\%) under IoT, proving that explicit
    constraint (via Anti-Purpose) prevents branching drift.
  \end{itemize}
\end{enumerate}

Beyond the formal framework and empirical lifecycle, we present a
longitudinal case study (Section 6) of SPAR (Structured
Persona-Argumentation for Reasoning), a multi-agent debate protocol that
independently evolved through 75 iterations under IoT governance,
validating the lifecycle's utility in a production system. The SPAR
protocol is open-source at https://github.com/synthanai/spar-kit and the
IoT framework, including all experimental data and prompts, is available
at https://github.com/synthanai/intent-of-thought.

We discuss our limitations honestly in Section 7: at current model
scales, single-path chain-of-thought remains an unreasonably robust
default, and the topology selection function's true empirical value lies
in catastrophic failure protection rather than theoretical optimisation.
The central thesis of this work is therefore not that IoT perfectly
predicts cognitive maps, but that \textbf{structured intent governance
improves aggregate reasoning resilience and drastically reduces
catastrophic topology mismatch, particularly for resource-constrained
models}. As we transition toward compounding agentic networks, we argue
that the locus of engineering value must shift from merely prompting for
answers toward architecting reliable \emph{intent structures} that
survive long-horizon inference.

The remainder of this paper is organised as follows. Section 2 surveys
related work on reasoning topologies, BDI agent governance, and failure
analysis. Section 3 formalises the IoT framework. Section 4 introduces
the operational lifecycle. Section 5 details our exhaustive experimental
results. Section 6 presents the longitudinal case study. Section 7
contextualises our findings within practical deployment scenarios, and
Section 8 concludes.

\section{Background and Related Work}\label{background-and-related-work}

This section surveys three research threads that converge on the problem
addressed by IoT: reasoning topology design, intent and governance in AI
systems, and reasoning failure analysis.

\subsection{Reasoning Topologies for Large Language
Models}\label{reasoning-topologies-for-large-language-models}

Chain-of-thought (CoT) prompting {[}Wei et al., 2022{]} demonstrated
that instructing language models to ``think step by step'' substantially
improves performance on arithmetic, commonsense, and symbolic reasoning
benchmarks. Subsequent work extended this linear structure.
Tree-of-thought (ToT) {[}Yao et al., 2023{]} introduces branching and
backtracking, enabling deliberate exploration of alternative reasoning
paths with the ability to prune unpromising branches. Graph-of-thought
(GoT) {[}Besta et al., 2024{]} generalises further, permitting arbitrary
merging and refinement across reasoning nodes, which supports non-linear
problem structures where partial solutions interact.
Algorithm-of-thought (AoT) {[}Hong et al., 2024{]} leverages algorithmic
patterns (e.g., depth-first search, dynamic programming) to impose
structured exploration within a single generation pass.

Recent advances have shifted from static topology definition to dynamic
routing and meta-reasoning. Framework of Thoughts (FoT) {[}Fricke et
al., 2026{]} provides a foundational framework for dynamically switching
between CoT, ToT, and GoT based on task complexity. Intention
Chain-of-Thought (ICoT) {[}Li et al., 2026{]} introduces dynamic routing
specifically for code generation, guided by intention states. Similarly,
Meta-Reasoner {[}Sui et al., 2026{]} offers dynamic guidance for
optimised inference-time reasoning by intervening during the generation
process. In multi-agent contexts, AMAS {[}Leong et al., 2025{]}
adaptively determines communication topologies between agents rather
than relying on fixed structures, while STELAR-VISION {[}Li et al.,
2026{]} applies self-topology-aware learning to efficient
vision-language tasks.

Each topology imposes structural commitments. CoT constrains reasoning
to a single sequential path. ToT introduces parallelism but requires an
evaluation function to select among branches. GoT permits cycles and
merges but requires careful state management. The choice of topology is
therefore not stylistic; it determines the reasoning architecture's
capabilities and failure modes.

Despite these advances in dynamic routing and meta-reasoning, the
topology governance problem remains largely open at the architectural
level. While frameworks like FoT and Meta-Reasoner optimise topology
dynamically during execution, they operate as inference-time optimisers
rather than pre-reasoning governance layers. No established mechanism
provides a strict, auditable bridge between \emph{why} an agent is
reasoning (its formal intent specification) and \emph{how} it structures
that reasoning, complete with a lifecycle for detecting and diagnosing
structural drift. This is the gap IoT addresses.

\subsection{Intent and Governance in AI
Systems}\label{intent-and-governance-in-ai-systems}

The notion that AI systems should be governed by explicit intent has
precedents across several domains, none of which address reasoning
topology governance directly.

\textbf{Prompt engineering and instruction tuning} {[}Brown et al.,
2020{]} provide mechanisms for specifying \emph{what} a model should
produce, but do not govern \emph{how} it should reason internally.
Similarly, \textbf{structured output formats} (e.g., JSON mode) only
constrain the final form, not the generative process.

Broader frameworks like \textbf{user intent modelling} in search
{[}Broder, 2002{]} and \textbf{AI safety tiers} (e.g., Anthropic's ASL)
operate either too broadly at the dialogue/system level or ignore
task-specific reasoning governance entirely. A model at ASL-2 uses the
same reasoning topology whether the task is routine or safety-critical.

IoT draws its primary inspiration from \textbf{BDI
(Belief-Desire-Intention) architectures} {[}Bratman, 1987; Cohen and
Levesque, 1990{]}, which formalise agent commitment. IoT adapts this
tradition for LLMs: reasoning persists until the Success Signal is
satisfied, the Purpose is judged unachievable, or Anti-Purpose is
violated.

The common gap across these precedents is that none ask: ``Given this
purpose, which reasoning topology should govern the process, and what
should we do when it fails?''

\subsection{Reasoning Failure
Analysis}\label{reasoning-failure-analysis}

When LLM reasoning fails, existing approaches diagnose failure at
different levels of abstraction, none of which address topology-level
diagnosis.

When LLMs fail, existing approaches diagnose the failure at the trace
level rather than the topology level. \textbf{Chain-of-thought error
analysis} {[}Ling et al., 2023{]} and \textbf{process reward models}
{[}Lightman et al., 2023{]} evaluate the correctness of individual
reasoning steps. However, step-correctness does not guarantee
topology-correctness: every step may be locally valid while the overall
structure is fundamentally misaligned with the intended task.

Self-corrective frameworks like \textbf{Reflexion} {[}Shinn et al.,
2023{]} improve outputs iteratively but reflect inherently on the
\emph{result}, not the \emph{topology selection}. An agent using
Reflexion may retry endlessly with the wrong structural approach,
producing refinements constrained by a flawed architectural commitment.

None of these approaches ask: ``Was the right reasoning topology
selected for this purpose?'' This is the exact question Retrospective
Judgement (Section 4.3) is designed to answer.

\subsection{Positioning: The Differentiation
Table}\label{positioning-the-differentiation-table}

Table 1 positions IoT relative to existing approaches across four
dimensions.

{\small\begin{longtable}[]{@{}
  >{\raggedright\arraybackslash}p{(\linewidth - 8\tabcolsep) * \real{0.1447}}
  >{\raggedright\arraybackslash}p{(\linewidth - 8\tabcolsep) * \real{0.2500}}
  >{\raggedright\arraybackslash}p{(\linewidth - 8\tabcolsep) * \real{0.1711}}
  >{\raggedright\arraybackslash}p{(\linewidth - 8\tabcolsep) * \real{0.2368}}
  >{\raggedright\arraybackslash}p{(\linewidth - 8\tabcolsep) * \real{0.1974}}@{}}
\toprule\noalign{}
Dimension
 & Prompt Engineering
 & IoT Capture
 & Failure Analysis
 & IoT Judgement
 \\
\midrule\noalign{}
\endhead
\bottomrule\noalign{}
\endlastfoot
\textbf{Input} & Task instructions & Governance triple (Purpose,
Anti-Purpose, Success Signal) & Reasoning trace & Trace + original IoT
spec \\
\textbf{Governs} & Task execution & Topology selection & Step
correctness & Topology + intent correctness \\
\textbf{Output} & Task result & Selected topology & Error location &
Diagnosis + corrective action \\
\textbf{Feedback} & None (open loop) & Fidelity gradient & Debug
information & Learning loop (improved capture) \\
\textbf{Operates at} & Instruction level & Governance level & Trace
level & Lifecycle level \\
\end{longtable}}

\emph{Table 1: Positioning of IoT relative to existing approaches. IoT
operates at the governance level (connecting intent to topology) and the
lifecycle level (connecting failure to learning). Existing approaches
operate at the instruction or trace level.}

\section{The IoT Framework}\label{the-iot-framework}

This section presents the core IoT governance layer: the intent triple
(Section 3.1), the topology selection function (Section 3.2), and the
drift detection protocol (Section 3.3). We then discuss the empirical
qualification that situates the selection function's value (Section
3.4).

\subsection{The Intent Triple}\label{the-intent-triple}

IoT specifies the \emph{intent} of a reasoning task through three
complementary primitives:

\textbf{Purpose.} The desired outcome of the reasoning process. Purpose
answers ``why are we reasoning?'' and specifies the end-state that the
agent seeks to achieve. Unlike a task description (which specifies
\emph{what} to do), Purpose specifies \emph{why} the reasoning matters
and what constitutes a good outcome. For instance, a Purpose might be
``identify the optimal investment allocation across three asset
classes'' rather than ``solve this optimisation problem.''

\textbf{Anti-Purpose.} The set of outcomes that would render the
reasoning worthless. Anti-Purpose answers ``what must we avoid?'' and
specifies explicit failure conditions. Anti-Purpose is operationally
distinct from constraints: where constraints limit the solution space,
Anti-Purpose identifies what would make the entire reasoning episode a
waste. For example: ``avoid recommending allocations that ignore tax
implications'' is a constraint, while ``reasoning that treats the three
asset classes as independent when they are correlated'' is an
Anti-Purpose. This distinction draws on the BDI tradition of commitment
termination {[}Cohen and Levesque, 1990{]}, where an agent abandons a
commitment not because the solution is infeasible but because the
reasoning has become purposeless.

\textbf{Success Signal.} The criteria by which the agent (or an
observer) determines that reasoning has achieved its Purpose. Success
Signal answers ``how will we know we succeeded?'' and provides
evaluation criteria that are checkable during and after reasoning.
Success Signals may be binary (achieved or not), graded (degree of
satisfaction), or multi-dimensional (a set of criteria).

These three primitives are complementary and non-redundant:

\begin{itemize}
\tightlist
\item
  Purpose without Anti-Purpose permits reasoning drift into technically
  valid but practically useless territory.
\item
  Purpose without Success Signal provides no mechanism for detecting
  when reasoning should terminate.
\item
  Anti-Purpose without Purpose is unconstrained negation: the system
  knows what to avoid but not what to achieve.
\end{itemize}

The triad functions as a pre-reasoning checkpoint: before selecting a
reasoning topology, the agent must specify all three components.

\textbf{The IoT Formula:}

\begin{quote}
\textbf{IoT = (Purpose, Anti-Purpose, Success Signal)}
\end{quote}

This triple is the input to the topology selection function defined
below.

\subsection{Topology Selection
Function}\label{topology-selection-function}

The selection function takes the IoT specification and the problem
context as input and produces a ranked list of recommended topologies
from the available set: Chain-of-Thought (CoT), Tree-of-Thought (ToT),
Graph-of-Thought (GoT), Algorithm-of-Thought (AoT), and Hybrid
combinations.

\textbf{Algorithm 1: Topology Selection}

\begin{verbatim}
Input:  IoT specification (Purpose, Anti-Purpose, Success Signal), problem context
Output: Ranked list of recommended topologies

Step 1. Extract purpose type.
  Analyse Purpose to identify the dominant reasoning requirement:
  sequential derivation, parallel exploration, interconnected
  analysis, hierarchical classification, or multi-phase complex.

Step 2. Match purpose type to topology.
  Use Table 2 to identify the topology whose structural
  properties best align with the identified purpose type.

Step 3. Apply Anti-Purpose constraints.
  Check whether the candidate topology would structurally
  violate Anti-Purpose. If so, demote it and promote the next candidate.

Output: Ordered list, most aligned topology first.
\end{verbatim}

\textbf{Table 2: IoT topology selection mapping.}

{\small\begin{longtable}[]{@{}
  >{\raggedright\arraybackslash}p{(\linewidth - 6\tabcolsep) * \real{0.1573}}
  >{\raggedright\arraybackslash}p{(\linewidth - 6\tabcolsep) * \real{0.2135}}
  >{\raggedright\arraybackslash}p{(\linewidth - 6\tabcolsep) * \real{0.2472}}
  >{\raggedright\arraybackslash}p{(\linewidth - 6\tabcolsep) * \real{0.3820}}@{}}
\toprule\noalign{}
Purpose Type
 & Structural Signal
 & Recommended Topology
 & Rationale
 \\
\midrule\noalign{}
\endhead
\bottomrule\noalign{}
\endlastfoot
Sequential derivation & Single valid path, step dependencies & CoT &
Linear chain; each step builds on the previous \\
Parallel exploration & Multiple viable approaches, uncertain best path &
ToT & Branching enables evaluation of alternatives \\
Interconnected analysis & Non-linear relationships, feedback loops & GoT
& Graph structure supports refinement and merging \\
Hierarchical classification & Problem type matters before details & AoT
& Abstraction first, then type-specific reasoning \\
Multi-phase complex & Requires multiple reasoning modes & Hybrid &
Sequential phases with different structural needs \\
\end{longtable}}

The mapping is not deterministic. Purpose determines the general
topology class, Anti-Purpose adds constraints (e.g., ``avoid premature
commitment to a single path'' favours ToT over CoT), and Success Signal
informs the evaluation mechanism within the selected topology.

\subsection{Intent Drift Detection}\label{intent-drift-detection}

Reasoning drift occurs when a multi-step reasoning process gradually
diverges from its stated purpose, producing outputs that are internally
coherent but misaligned with the original intent. IoT addresses drift at
the topology-governance level through continuous monitoring.

\textbf{Drift Score.} The drift detection function measures the
alignment between the current reasoning trajectory and the stated
Purpose, producing a score between 0 (fully aligned) and 1 (completely
diverged). When this drift score exceeds a predefined threshold, the
system triggers corrective action.

\textbf{Algorithm 2: Intent Drift Detection}

\begin{verbatim}
Input:  IoT specification (Purpose, Anti-Purpose, Success Signal), reasoning trace
Output: Continue, correct, switch topology, or terminate

Step 1. Anchor. Record Purpose as the reference point.

Step 2. Monitor. At regular intervals, compute drift score.

Step 3. Drift check. If drift score exceeds threshold:
  a. Re-read the IoT specification.
  b. Evaluate whether the current topology is still
     appropriate for Purpose.
  c. If topology is mismatched: re-run Algorithm 1 with
     updated context.
  d. If topology is appropriate: add a purpose-realignment
     step to the reasoning chain.

Step 4. Anti-Purpose violation. If reasoning enters territory
  identified by Anti-Purpose: trigger immediate correction.

Step 5. Completion. If Success Signal is satisfied: terminate.
  If Purpose is believed unachievable: terminate and report.
\end{verbatim}

This protocol draws on BDI commitment termination {[}Cohen and Levesque,
1990{]}: reasoning persists until the Purpose is achieved, believed
unachievable, or superseded by Anti-Purpose violation.

\subsection{Empirical Qualification}\label{empirical-qualification}

The topology selection mapping in Table 2 was designed prescriptively:
``match topology to task type.'' Our experimental evidence (Section 5)
reveals that this prescription must be qualified.

At current model scales (7B to 120B parameters), chain-of-thought (CoT)
is empirically robust across all task types, including those Table 2
assigns to GoT or ToT. CoT scored 2.47 on interconnected tasks where GoT
scored 2.20 (N = 131, Section 5.4). The selection function's empirical
value is therefore not in \emph{optimal} topology recommendation but in
\emph{failure protection}: preventing the catastrophic quality loss that
occurs when the wrong topology is selected without governance (e.g., ToT
on interconnected tasks: 1.00/3.00).

This qualification does not invalidate the selection function. It
situates its contribution accurately: governance provides a safety net
that makes topology misselection survivable, and an uplift mechanism
that improves all topologies. As models scale beyond the ranges tested
here, the selection function's prescriptive role may become empirically
necessary. At current scales, its value is architectural preparedness
combined with measurable quality improvement.

\section{The IoT Lifecycle}\label{the-iot-lifecycle}

The framework presented in Section 3 operates under three simplifying
assumptions: (i) the intent triple is \emph{given} by a competent
specifier, (ii) the triple \emph{correctly} captures the actual
reasoning need, and (iii) governance operates \emph{forward only}, from
specification through selection to monitoring. This section relaxes all
three assumptions by introducing the operational lifecycle that
completes the IoT governance layer.

\subsection{The Capture Spectrum}\label{the-capture-spectrum}

Intent elicitation for reasoning governance can be organised as a
five-mode hierarchy that produces IoT triples of varying fidelity. We
term this the \emph{Capture Spectrum} and frame it as an intent
elicitation maturity hierarchy.

\textbf{Table 3: The Capture Spectrum.}

{\small\begin{longtable}[]{@{}
  >{\raggedright\arraybackslash}p{(\linewidth - 8\tabcolsep) * \real{0.0893}}
  >{\raggedright\arraybackslash}p{(\linewidth - 8\tabcolsep) * \real{0.0476}}
  >{\raggedright\arraybackslash}p{(\linewidth - 8\tabcolsep) * \real{0.2679}}
  >{\raggedright\arraybackslash}p{(\linewidth - 8\tabcolsep) * \real{0.2679}}
  >{\raggedright\arraybackslash}p{(\linewidth - 8\tabcolsep) * \real{0.3274}}@{}}
\toprule\noalign{}
Mode
 & Level
 & Mechanism
 & Triple Fidelity
 & Example
 \\
\midrule\noalign{}
\endhead
\bottomrule\noalign{}
\endlastfoot
Zero & L0 & System infers intent from task context alone & Low: Purpose
inferred, Anti-Purpose absent, Success Signal generic & ``Solve this
problem'' leads to system defaulting to CoT \\
Implicit & L1 & Lightweight extraction from user phrasing and domain &
Medium-low: Purpose extracted, Anti-Purpose partial & ``Compare these
options carefully'' carries an implicit ToT signal \\
Prompted & L2 & Structured questions elicit Purpose, Anti-Purpose,
Success Signal explicitly & Medium-high: full triple, user-specified &
System asks: ``What must this analysis avoid?'' \\
Collaborative & L3 & Interactive refinement loop: draft, feedback,
revise & High: validated triple via dialogue & Multi-turn IoT
specification with user correction \\
Learned & L4 & Model has internalised intent-topology mappings from
training & Variable: training-dependent & Model auto-generates IoT
triple from task embedding \\
\end{longtable}}

The hierarchy is not purely ordinal: ``higher is always better'' does
not hold. L0 is appropriate for low-stakes, familiar tasks; L3
introduces interaction cost that may be unwarranted for routine queries.
The spectrum describes \emph{available capability}, not prescribed
usage. A mature system selects its capture mode based on task
criticality, user context, and confidence.

\textbf{How capture confidence drives topology selection.} The system
assigns a confidence score (0 to 1) to each captured triple. When this
confidence is low, the system either escalates its capture mode (e.g.,
from L0 to L2) or defaults to a safe, general-purpose topology. When
confidence is high, the system can recommend specialised topologies
because the governance triple is precise enough to justify the
structural commitment.

\textbf{Architectural distinction from prompt engineering.} The
strongest objection to the Capture Spectrum is that it constitutes
``prompt engineering with extra steps.'' The decisive rebuttal is that
the capture modes imply different \emph{control surfaces}, not merely
different textual instructions:

\begin{enumerate}
\def\labelenumi{\arabic{enumi}.}
\tightlist
\item
  \textbf{Typed artefacts}: intent is materialised as a structured
  contract, not free-form text.
\item
  \textbf{Policies}: the system follows explicit rules for when to ask
  versus when to act, governed by fidelity thresholds.
\item
  \textbf{State machines}: capture proceeds through defined phases
  (draft, clarify, commit) with auditable transitions.
\item
  \textbf{Telemetry}: field-level missingness, correction rates, and
  capture mode distributions are measurable.
\end{enumerate}

A prompt produces a task output consumed by the user; an IoT Capture
produces a governance triple consumed by the topology selection
function. Different type signatures entail different architectural
roles.

\subsection{Capture-Topology Confidence
Interaction}\label{capture-topology-confidence-interaction}

The confidence score creates a correspondence between capture quality
and topology specialisation:

\begin{itemize}
\tightlist
\item
  \textbf{Low-fidelity capture} (L0 to L1): the selection function
  should prefer robust, general-purpose topologies. Our experimental
  evidence (Section 5) confirms that chain-of-thought serves as a safe
  default at these fidelity levels, because CoT is empirically robust
  across task types.
\item
  \textbf{High-fidelity capture} (L3 to L4): the selection function can
  recommend specialised topologies (GoT, hybrid strategies), because the
  governance triple is precise enough to justify the structural
  commitment and mitigate the risk of misselection.
\end{itemize}

This interaction implies that topology selection is not merely a
function of task content but of \emph{how well the system understands
what it is reasoning about}.

\subsection{Retrospective Judgement}\label{retrospective-judgement}

When reasoning fails despite governance, the system must diagnose
\emph{why}. Retrospective Judgement is backward intent reconstruction:
given a failed reasoning trace and the original IoT specification,
reconstruct what the intent \emph{should have been} and classify the
failure mode.

\textbf{How Retrospective Judgement works.} When a reasoning episode
fails, the judgement function takes the failed trace and the original
IoT triple as inputs, and produces three outputs: the diagnosed failure
mode, a reconstructed version of what the triple should have been, and a
recommended corrective action.

The three failure modes form a diagnostic hierarchy:

\textbf{False Capture.} The IoT triple was incorrectly specified. The
Purpose does not match the actual reasoning need. Example: a clinical
triage task specifies Purpose as ``identify the most likely diagnosis''
when the actual need is ``identify the most dangerous diagnosis that
must be ruled out.'' The topology selection was correct \emph{given} the
stated Purpose, but the stated Purpose was wrong.

\textbf{False Selection.} The IoT triple was correct, but the selection
function chose the wrong topology. Example: an interconnected analysis
task correctly specifies Purpose as ``identify interactions between
factors,'' but the system selects CoT instead of GoT, producing a linear
analysis that misses cross-factor relationships.

\textbf{False Execution.} The IoT triple and topology selection were
both correct, but reasoning drifted during execution. The drift
detection either failed to trigger or triggered too late.

\textbf{Diagnostic algorithm.} Judgement proceeds by comparing the
failed trace against the original IoT specification:

\begin{verbatim}
Algorithm 3: Retrospective Judgement

Input:  Failed trace, selected topology,
        original IoT = (Purpose, Anti-Purpose, Success Signal)
Output: Failure mode, reconstructed IoT, corrective action

Step 1. Reconstruct intent from trace.
  Analyse the trace to infer what purpose the reasoning was
  actually pursuing. Call this the "actual purpose."

Step 2. Compare actual purpose with stated Purpose.
  If they differ: diagnose False Capture.
    Action: Escalate capture mode (e.g., L0 to L2).
  If they match: proceed to Step 3.

Step 3. Evaluate topology-purpose alignment.
  Using Table 2, check whether the selected topology was appropriate for Purpose.
  If mismatched: diagnose False Selection.
    Action: Re-run Algorithm 1 with confirmed Purpose.
  If appropriate: diagnose False Execution.
    Action: Lower drift threshold, add checkpoints.

Output: (failure_mode, reconstructed_IoT, corrective_action)
\end{verbatim}

\subsection{The Learning Loop}\label{the-learning-loop}

Capture and Judgement are structurally symmetric: Capture moves from
intent to topology (forward), while Judgement moves from failed topology
to reconstructed intent (backward). Both operate on the same IoT triple
but from opposite temporal positions.

The Learning Loop connects Judgement outputs to improved Capture:

\begin{enumerate}
\def\labelenumi{\arabic{enumi}.}
\tightlist
\item
  \textbf{False Capture corrections} feed back into capture policies,
  updating the system's thresholds for when to escalate from L0 to L2 or
  L3.
\item
  \textbf{False Selection corrections} feed back into the selection
  function, refining the topology-purpose mappings in Table 2 based on
  observed failures.
\item
  \textbf{False Execution corrections} feed back into drift detection,
  adjusting the threshold and checkpoint frequency.
\end{enumerate}

This feedback loop closes the lifecycle:

\begin{quote}
\textbf{Capture \textgreater{} Select \textgreater{} Monitor
\textgreater{} Judge \textgreater{} Learn \textgreater{} Capture}
\end{quote}

The loop is not merely circular; it is adaptive. Each iteration improves
the system's ability to capture intent, select topologies, and detect
failures. The convergence condition is when Judgement produces no new
corrections, a state we term \emph{governance equilibrium}.

\subsection{L4 Learned Capture: A Future
Direction}\label{l4-learned-capture-a-future-direction}

L4 is the most architecturally distinctive capture mode and the most
speculative. Where L0 through L3 operate at inference time, L4 operates
at training time: the model has internalised intent-topology mappings
from prior data, enabling auto-generation of IoT triples without
explicit elicitation.

Training-time tool integration provides a partial precedent: Toolformer
{[}Schick et al., 2023{]} trains models to decide \emph{when and how} to
invoke external tools. L4 extends this: instead of learning to produce
\emph{tool invocations}, the model learns to produce \emph{governance
triples}. We do not evaluate L4 empirically in this paper and include it
as a theoretical contribution that defines the spectrum's upper bound.

With the theoretical lifecycle established, from explicit L0 capture up
to speculative L4 internalisation, the architectural question naturally
shifts from \emph{how} to govern topologies to \emph{whether} this
governance actually yields measurable improvements. Section 5
transitions from this theoretical foundation into the empirical
evidence, explicitly tracking whether forcing these structural
constraints successfully prevents the catastrophic cognitive drift
predicted by our failure taxonomy.

\section{Evidence}\label{evidence}

We evaluate the IoT governance layer through three experiments designed
to test its core claims: that governance improves reasoning quality
(Experiment 1), that the Anti-Purpose primitive contributes measurably
(Experiment 5), and that topology misselection produces predictable
degradation patterns (Experiment 8). A fourth experiment tests
generalisation to frontier API models (Experiment 9).

\subsection{Experimental Design}\label{experimental-design}

\textbf{Models.} We evaluate 7 models spanning consumer-grade (7B
parameters) to large mixture-of-experts (120B parameters, 12B active):

{\small\begin{longtable}[]{@{}lccc@{}}
\toprule\noalign{}
Model & Parameters & Architecture & Backend \\
\midrule\noalign{}
\endhead
\bottomrule\noalign{}
\endlastfoot
qwen2.5:7b & 7B & Dense & Local (Ollama) \\
mistral:latest & 7B & Dense & Local (Ollama) \\
deepseek-r1:7b & 7B & Dense & Local (Ollama) \\
qwen3.5:9b & 9B & Dense & Local (Ollama) \\
gemma3:12b & 12B & Dense & Local (Ollama) \\
phi4:14b & 14B & Dense & Local (Ollama) \\
nvidia/nemotron-120b & 120B (12B active) & MoE & Cloud (OpenRouter) \\
\end{longtable}}

One model (qwen3.5:9b) was excluded from analysis due to insufficient
data: 90\% of its responses were empty (generation failures), yielding
only 1 usable response for the baseline condition. One cloud model
(z-ai/glm-5-turbo) was included in the experiment but showed
non-significant results (-8.3\%, per-task analysis revealed inconsistent
directionality; see Section 7). Results reported below are for the 7
models with sufficient data.

\textbf{Tasks.} 18 reasoning tasks across 6 domains (clinical diagnosis,
software debugging, mathematical proof, logistics optimisation,
architecture comparison, algorithm tracing). Tasks are classified by
optimal reasoning structure: \emph{sequential} (single-path derivation),
\emph{parallel} (multiple viable approaches), and \emph{interconnected}
(non-linear relationships requiring cross-factor analysis).

\textbf{Topologies.} Four topologies tested: chain-of-thought (CoT),
tree-of-thought (ToT), graph-of-thought (GoT), and a baseline condition
(no topology instruction).

\textbf{Conditions.} Two governance conditions:

\begin{itemize}
\tightlist
\item
  \emph{Baseline}: task prompt only, no IoT specification.
\item
  \emph{IoT L2}: task prompt preceded by an explicit IoT governance
  triple specifying Purpose, Anti-Purpose, and Success Signal (Prompted
  capture, Table 3).
\end{itemize}

\textbf{Scoring.} All responses scored on a 0 to 3 scale (0: wrong, 1:
partial, 2: correct, 3: excellent) using an LLM-as-judge approach
{[}Zheng et al., 2023{]} with phi4:14b as evaluator. Each response was
scored against a domain-expert-authored ground truth rubric. Judge
self-consistency was measured at 98\% (91/93 re-scored responses
matched).

\textbf{Reproducibility.} All local models were run via Ollama with
temperature 0.7 and default sampling parameters. Cloud models were
accessed via OpenRouter API with temperature 0.7. No random seeds were
fixed (each response represents a single stochastic sample). Prompt
templates for all three experiments, including the IoT governance triple
format and the LLM-as-judge rubric, are provided in the supplementary
materials.

\textbf{Infrastructure.} Local models served via Ollama; cloud models
accessed via OpenRouter API. Total scored evaluations: 705 across all
experiments.

\subsection{Experiment 1: Governance Impact (N =
486)}\label{experiment-1-governance-impact-n-486}

\textbf{Hypothesis.} Adding an IoT governance triple (L2 Prompted
capture) before reasoning improves output quality relative to an
ungoverned baseline.

\subsubsection{Results}\label{results}

\textbf{Table 4: Overall governance impact.}

{\small\begin{longtable}[]{@{}lcc@{}}
\toprule\noalign{}
Condition & Mean Score (0 to 3) & N \\
\midrule\noalign{}
\endhead
\bottomrule\noalign{}
\endlastfoot
Baseline (no IoT) & 2.10 & 240 \\
IoT L2 (governance triple) & \textbf{2.43} & 246 \\
\textbf{Uplift} & \textbf{+0.33 (+15.7\%)} & \\
\end{longtable}}

The governance uplift is consistent across the majority of models. A
paired t-test on the subset of exactly matched reasoning episodes (N =
209 matched pairs) confirms this overall uplift is statistically
significant (p \textless{} 0.001). Table 5 reports the granular
per-model breakdown.

\textbf{Table 5: Per-model governance impact (sorted by uplift).}

{\small\begin{longtable}[]{@{}lccccc@{}}
\toprule\noalign{}
Model & Params & Baseline & IoT L2 & Change & Uplift \\
\midrule\noalign{}
\endhead
\bottomrule\noalign{}
\endlastfoot
qwen2.5:7b & 7B & 2.08 & 2.77 & +0.69 & \textbf{+33.3\%} \\
mistral:latest & 7B & 1.62 & 2.10 & +0.49 & \textbf{+30.2\%} \\
gemma3:12b & 12B & 2.13 & 2.61 & +0.47 & \textbf{+22.2\%} \\
nemotron-120b & 120B & 2.27 & 2.64 & +0.37 & +16.3\% \\
phi4:14b & 14B & 2.34 & 2.71 & +0.37 & +15.7\% \\
deepseek-r1:7b & 7B & 2.09 & 2.34 & +0.25 & +12.0\% \\
z-ai/glm-5-turbo & \textasciitilde14B & 2.31 & 2.12 & -0.19 & -8.3\% \\
\end{longtable}}

\subsubsection{Finding 1: Governance as Cognitive
Equaliser}\label{finding-1-governance-as-cognitive-equaliser}

The governance uplift is inversely correlated with model capability:
smaller models (7B) gain +30 to 33\%, while larger models (14B to 120B)
gain +12 to 16\%. This pattern suggests that IoT governance functions as
a \emph{cognitive equaliser}: it compensates for weaker models'
inability to self-structure reasoning, providing external architectural
scaffolding that stronger models partially possess internally.

\subsubsection{Finding 2: Governance as Topology Safety
Net}\label{finding-2-governance-as-topology-safety-net}

Table 6 reports the per-topology governance impact.

\textbf{Table 6: Per-topology governance impact (sorted by uplift).}

{\small\begin{longtable}[]{@{}lcccc@{}}
\toprule\noalign{}
Topology & Baseline & IoT L2 & Change & Uplift \\
\midrule\noalign{}
\endhead
\bottomrule\noalign{}
\endlastfoot
ToT & 1.67 & 2.28 & +0.62 & \textbf{+37.0\%} \\
GoT & 1.96 & 2.53 & +0.57 & \textbf{+29.1\%} \\
Baseline (none) & 2.37 & 2.48 & +0.12 & +4.9\% \\
CoT & 2.35 & 2.44 & +0.09 & +3.8\% \\
\end{longtable}}

Governance rescues the worst-performing topologies: ToT gains +37\% and
GoT gains +29\%. CoT, already robust, gains only +3.8\%. The governance
benefit scales inversely with topology quality: the more a topology
needs structural guidance, the more governance helps.

This finding has an important practical implication: IoT governance is
not merely a quality improvement mechanism; it is a \emph{safety net}
that makes topology misselection survivable. Without governance,
selecting the wrong topology produces catastrophic results (see
Experiment 8 below). With governance, the same misselection produces
usable, if suboptimal, results.

\subsection{Experiment 5: Anti-Purpose Ablation (N =
88)}\label{experiment-5-anti-purpose-ablation-n-88}

\textbf{Hypothesis.} The Anti-Purpose primitive contributes measurably
to reasoning quality beyond what Purpose and Success Signal provide.

\subsubsection{Design}\label{design}

Three conditions, each using L2 Prompted capture with one component
ablated: - \emph{Full Triple}: complete IoT specification (Purpose,
Anti-Purpose, Success Signal). - \emph{No Anti-Purpose}: Purpose and
Success Signal only. - \emph{No Success Signal}: Purpose and
Anti-Purpose only.

\subsubsection{Results}\label{results-1}

\textbf{Table 7: Anti-Purpose ablation.}

{\small\begin{longtable}[]{@{}
  >{\raggedright\arraybackslash}p{(\linewidth - 8\tabcolsep) * \real{0.2157}}
  >{\centering\arraybackslash}p{(\linewidth - 8\tabcolsep) * \real{0.2157}}
  >{\centering\arraybackslash}p{(\linewidth - 8\tabcolsep) * \real{0.0980}}
  >{\centering\arraybackslash}p{(\linewidth - 8\tabcolsep) * \real{0.0980}}
  >{\centering\arraybackslash}p{(\linewidth - 8\tabcolsep) * \real{0.3725}}@{}}
\toprule\noalign{}
Condition
 & Mean Score
 & Std Dev
 & N
 & Change vs Full
 \\
\midrule\noalign{}
\endhead
\bottomrule\noalign{}
\endlastfoot
Full Triple (Purpose, Anti-Purpose, Success Signal) & \textbf{2.20} &
0.81 & 30 & Ref. \\
No Success Signal (Purpose, Anti-Purpose) & 2.17 & 0.83 & 30 & -0.03 \\
No Anti-Purpose (Purpose, Success Signal) & 2.11 & 0.88 & 28 & -0.09 \\
\end{longtable}}

\subsubsection{Finding 3: Anti-Purpose as Drift
Guard}\label{finding-3-anti-purpose-as-drift-guard}

Removing Anti-Purpose produces the lowest scores, confirming that
explicit negative guardrails contribute to reasoning quality. We note,
however, the high variance (standard deviation around 0.85 on a 3-point
scale) and the modesty of the effect size (+4\% uplift when Anti-Purpose
is present). A paired t-test indicates this difference is directionally
consistent but not statistically significant (p = 0.14) at current
sample sizes.

Rather than a definitive proof of necessity, we interpret this finding
as consistent with the view that Anti-Purpose functions as a drift
guard: it operates as a boundary constraint rather than a primary driver
of reasoning success. Its absence does not cause catastrophic failure;
it permits gradual drift that accumulates over multi-step reasoning. We
expect the effect to be larger on longer reasoning chains and more
open-ended generative tasks, remaining a structural hypothesis for
future work.

\subsection{Experiment 8: The Confusion Matrix (N =
131)}\label{experiment-8-the-confusion-matrix-n-131}

\textbf{Hypothesis.} Using the wrong topology for a given task category
produces measurably worse outcomes, and the degradation pattern is
predictable from the topology-purpose mapping (Table 2).

\subsubsection{Design}\label{design-1}

Each of 3 task categories (sequential, parallel, interconnected) is
paired with each of 3 topologies (CoT, ToT, GoT), producing a 3x3
confusion matrix. Tasks are run across 5 models under IoT L2 governance.

\subsubsection{Results}\label{results-2}

\textbf{Table 8: The Confusion Matrix (Task Category x Topology).}

{\small\begin{longtable}[]{@{}lcccc@{}}
\toprule\noalign{}
Task Category & CoT & GoT & ToT & Optimal (Table 2) \\
\midrule\noalign{}
\endhead
\bottomrule\noalign{}
\endlastfoot
Sequential & \textbf{2.13} & 1.71 & 1.64 & CoT \\
Parallel & \textbf{2.07} & 1.27 & 2.00 & ToT \\
Interconnected & \textbf{2.47} & 2.20 & 1.00 & GoT \\
\end{longtable}}

\subsubsection{Finding 4: CoT Unreasonable
Robustness}\label{finding-4-cot-unreasonable-robustness}

CoT achieves the highest or near-highest score in \emph{every} task
category, including those Table 2 assigns to GoT (interconnected: 2.47
vs.~2.20) and ToT (parallel: 2.07 vs.~2.00). At the model scales tested
(7B to 14B), CoT is a universally safe default.

\subsubsection{Finding 5: Topology Misselection is
Catastrophic}\label{finding-5-topology-misselection-is-catastrophic}

The confusion matrix confirms that wrong topology selection can be
catastrophic. ToT on interconnected tasks scores 1.00, a complete
failure. GoT on parallel tasks scores 1.27. These are not gradual
degradations; they are structural mismatches that produce qualitatively
different (worse) outputs.

\subsubsection{Finding 6: The Combined
Insight}\label{finding-6-the-combined-insight}

Findings 4 and 5 together produce the paper's most practically important
result: \emph{you probably don't need to select a topology (CoT is
safe), but if you do, getting it wrong is catastrophic, and governance
makes the wrong selection survivable (Experiment 1, Table 6).}

This reframes the IoT framework's contribution: from ``always select the
optimal topology'' to ``default to CoT, escalate only when justified,
and use governance to protect against the consequences of topology
failure.''

\textbf{Table 9: Per-model confusion matrix scores.}

{\small\begin{longtable}[]{@{}lccc@{}}
\toprule\noalign{}
Model & CoT & GoT & ToT \\
\midrule\noalign{}
\endhead
\bottomrule\noalign{}
\endlastfoot
phi4:14b & 2.56 & 1.89 & 2.11 \\
gemma3:12b & 2.56 & 1.56 & 1.89 \\
qwen2.5:7b & 2.33 & 2.11 & 1.33 \\
deepseek-r1:7b & 2.22 & 1.88 & 1.33 \\
mistral:latest & 1.44 & 1.22 & 1.00 \\
\end{longtable}}

The CoT robustness pattern holds across all models tested, including the
weakest (mistral, 1.44). The cross-model consistency strengthens the
finding's generalisability within the tested parameter range.

\subsection{Experiment 9: Frontier Model Generalisation (N =
72)}\label{experiment-9-frontier-model-generalisation-n-72}

\textbf{Hypothesis.} The governance uplift observed in local 7B to 14B
models generalises to frontier API models, though the absolute magnitude
of the uplift will scale inversely with the model's baseline capability.

\subsubsection{Design}\label{design-2}

Two frontier models, gpt-4o-mini (\textasciitilde8B distilled) and
claude-3.5-haiku (\textasciitilde20B), were evaluated across the 18-task
comprehensive benchmark. Each model was tested in two conditions: an
unguided baseline and the full IoT L2 (Purpose, Anti-Purpose, Success
Signal) governance prompt. Responses were independently scored by
phi4:14b on the standard 0 to 3 rubric.

\subsubsection{Results}\label{results-3}

\textbf{Table 10: Frontier model governance impact.}

{\small\begin{longtable}[]{@{}lcccc@{}}
\toprule\noalign{}
Model & Baseline Mean & IoT L2 Mean & Change (Uplift) & N \\
\midrule\noalign{}
\endhead
\bottomrule\noalign{}
\endlastfoot
claude-3.5-haiku & 2.83 & 2.78 & -0.06 (-2.0\%) & 18 \\
gpt-4o-mini & 2.28 & 2.56 & +0.28 (+12.2\%) & 18 \\
\textbf{Overall Frontier} & 2.56 & 2.67 & \textbf{+0.11 (+4.3\%)} &
36 \\
\end{longtable}}

\subsubsection{Finding 7: The Capability Ceiling
Effect}\label{finding-7-the-capability-ceiling-effect}

The results confirm that governance uplift does not scale linearly with
model intelligence; it scales with \emph{model need}.

Claude 3.5 Haiku achieved a near-perfect baseline score (2.83/3.00). At
this level of baseline competence, explicit structural governance is
redundant and imposes a slight cognitive overhead, resulting in a
marginal loss (-2.0\%). The model already intrinsically executes the
guardrails that IoT makes explicit.

In contrast, gpt-4o-mini (baseline 2.28) demonstrated a significant
+12.2\% uplift when governed by the IoT triple, elevating its reasoning
reliability toward Haiku's baseline. This validates that the Intent of
Thought framework is highly effective for ``punching up'' the
reliability of smaller, faster, or distilled models, transforming
cost-effective inference engines into robust reasoning agents.

\section{Case Study: SPAR, A Multi-Agent Debate Protocol as
Self-Realised
IoT}\label{case-study-spar-a-multi-agent-debate-protocol-as-self-realised-iot}

The evidence in Section 5 evaluates IoT governance in controlled,
single-episode reasoning tasks. This section presents a longitudinal
case study of IoT governance operating in a production system: SPAR
(Structured Persona-Argumentation for Reasoning), a multi-agent debate
protocol that evolved through five distinct phases over 75 debates
spanning twelve months. The SPAR protocol and its implementation are
publicly available at https://github.com/synthanai/spar-kit.

Multi-agent debate, where multiple LLMs argue from different
perspectives to improve reasoning quality, has shown promise for fact
verification {[}Du et al., 2023{]}, mathematical reasoning {[}Liang et
al., 2024{]}, and complex decision-making {[}Chan et al., 2024{]}. The
SPAR case study is notable because IoT governance was not imposed
retroactively; the protocol \emph{independently} developed IoT-like
governance structures, which we later recognised as a self-realised
instance of the lifecycle described in Section 4.

\subsection{Context: The SPAR Protocol}\label{context-the-spar-protocol}

SPAR orchestrates structured debates between 4 to 8 LLM agents, each
assigned a distinct persona and argumentative stance, to produce
reasoned decisions on complex questions (e.g., technology selection,
strategic priorities, publication strategy). Each debate follows a
multi-round structure: agents state positions, challenge each other,
synthesise tensions, and converge toward a verdict with explicit dissent
recording.

The protocol was initially designed as a flat structure: agents debate,
a moderator synthesises. No governance layer existed. Over 75 debates,
the protocol evolved iteratively, with each modification motivated by
observed failure modes.

\subsection{The Five Phases of
Evolution}\label{the-five-phases-of-evolution}

\textbf{Phase 1: Role-based debate (SPAR v1.0 to v2.0, debates 1 to
12).} Each agent was assigned a role (Advocate, Skeptic, Theorist) and a
stance. The only governance was the moderator's prompt. Failure mode
observed: agents drifted into tangential arguments that were internally
coherent but purposeless. There was no mechanism to detect this drift or
to specify what the debate should avoid.

\emph{IoT interpretation}: This is L0 capture (zero mode). No intent
triple exists. The system infers purpose from context alone. Drift
detection is absent. This pattern is consistent with the ``naive
debate'' condition in {[}Du et al., 2023{]}, where agents debate without
structured governance.

\textbf{Phase 2: Purpose-driven debate (SPAR v2.1 to v3.0, debates 13 to
28).} After repeated drift, a Purpose statement was added to the debate
prompt: ``The goal of this debate is {[}X{]}.'' This reduced drift but
introduced a new failure mode: debates would satisfy the stated Purpose
through superficial consensus, avoiding genuine tension.

\emph{IoT interpretation}: Partial L2 capture. Purpose is explicit, but
Anti-Purpose and Success Signal are absent. Without Anti-Purpose, there
is no mechanism to prevent purposeless agreement.

\textbf{Phase 3: Anti-Purpose and Success Signal (SPAR v3.1 to v5.0,
debates 29 to 50).} Two additions solved the superficial consensus
problem. First, Anti-Purpose: ``This debate must NOT {[}Y{]}'' (e.g.,
``must not reach consensus without addressing the strongest
counterargument''). Second, Success Signal: explicit criteria for what
constitutes a valid debate output (e.g., ``must include dissent record,
confidence score, and review date'').

\emph{IoT interpretation}: Full L2 capture. The complete IoT triple
(Purpose, Anti-Purpose, Success Signal) governs each debate. This is the
configuration tested in Experiment 1 (Section 5.2), where it produced a
15.7\% quality uplift.

\textbf{Phase 4: Topology-aware debate (SPAR v5.1 to v6.0, debates 51 to
65).} With the governance triple established, the protocol began
adapting its \emph{reasoning structure} based on the debate's purpose.
For tradeoff decisions, agents were instructed to follow tree-of-thought
(branch and compare). For interconnected analyses, graph-of-thought (map
factors and relationships). The topology selection was governed by the
IoT triple: the Purpose determined the reasoning structure, the
Anti-Purpose constrained it, and the Success Signal validated the
output.

\emph{IoT interpretation}: The topology selection function (Algorithm 1)
is now operational. Each debate's structure is chosen based on its
governance specification, not by default.

\textbf{Phase 5: Retrospective judgement and learning (SPAR v6.1 to
v7.4, debates 66 to 75).} The protocol added systematic retrospection:
after each debate, a structured reflection document (a LOON, or
backward-looking reasoning trace) recorded what went well, what failed,
and what should change. Failed debates were diagnosed: Was the Purpose
wrong (False Capture)? Was the topology mismatched (False Selection)?
Did agents drift despite correct setup (False Execution)? These
diagnoses fed back into improved debate templates.

\emph{IoT interpretation}: The full lifecycle is operational. Capture
\textgreater{} Select \textgreater{} Monitor \textgreater{} Judge
\textgreater{} Learn \textgreater{} Capture. The Learning Loop closed
naturally over 75 iterations.

\subsection{Quantitative Evolution}\label{quantitative-evolution}

Table 11 summarises the structural progression across SPAR versions.

{\small\begin{longtable}[]{@{}
  >{\centering\arraybackslash}p{(\linewidth - 10\tabcolsep) * \real{0.0753}}
  >{\centering\arraybackslash}p{(\linewidth - 10\tabcolsep) * \real{0.1505}}
  >{\centering\arraybackslash}p{(\linewidth - 10\tabcolsep) * \real{0.0968}}
  >{\centering\arraybackslash}p{(\linewidth - 10\tabcolsep) * \real{0.2043}}
  >{\centering\arraybackslash}p{(\linewidth - 10\tabcolsep) * \real{0.2151}}
  >{\centering\arraybackslash}p{(\linewidth - 10\tabcolsep) * \real{0.2581}}@{}}
\toprule\noalign{}
Phase
 & SPAR Version
 & Debates
 & IoT Capture Level
 & Components Present
 & Failure Mode Addressed
 \\
\midrule\noalign{}
\endhead
\bottomrule\noalign{}
\endlastfoot
1 & v1.0 to v2.0 & 1 to 12 & L0 (Zero) & None & (none) \\
2 & v2.1 to v3.0 & 13 to 28 & Partial L2 & Purpose only & Drift \\
3 & v3.1 to v5.0 & 29 to 50 & L2 (Prompted) & Purpose, Anti-Purpose,
Success Signal & Superficial consensus \\
4 & v5.1 to v6.0 & 51 to 65 & L2 + Selection & Full triple + topology
function & Topology mismatch \\
5 & v6.1 to v7.4 & 66 to 75 & L2 + Lifecycle & Full lifecycle &
Recurrence of known failures \\
\end{longtable}}

\subsection{What the Case Study
Demonstrates}\label{what-the-case-study-demonstrates}

This case study provides three forms of evidence that complement the
controlled experiments in Section 5:

\begin{enumerate}
\def\labelenumi{\arabic{enumi}.}
\item
  \textbf{Organic emergence}: The IoT lifecycle was not designed
  top-down and applied to SPAR. It emerged bottom-up: each governance
  component was added in response to an observed failure mode. This
  suggests the lifecycle addresses real, recurrent needs rather than
  theoretical constructs.
\item
  \textbf{Longitudinal validation}: The controlled experiments (Section
  5) measure IoT governance in single episodes. The SPAR case study
  demonstrates governance improving \emph{across} episodes through the
  Learning Loop, validating the closed-loop lifecycle over sustained
  use.
\item
  \textbf{Anti-Purpose as the critical inflection point}: The transition
  from Phase 2 (Purpose only) to Phase 3 (full triple) produced the most
  noticeable qualitative improvement. This corroborates Experiment 5's
  finding that Anti-Purpose contributes measurably, and suggests the
  effect may compound over time as Anti-Purpose accumulates
  institutional knowledge of failure patterns.
\end{enumerate}

The case study does not claim controlled causality: other factors
(prompt refinement, model upgrades, user learning) contributed to the
protocol's improvement. The claim is structural: the IoT lifecycle
accurately describes the governance evolution that occurred, and the
protocol's observed failure modes map cleanly onto the three diagnostic
categories (False Capture, False Selection, False Execution) defined in
Section 4.3. The full SPAR protocol, including debate templates, persona
definitions, and governance configurations, is open-source at
https://github.com/synthanai/spar-kit.

\section{Discussion and Limitations}\label{discussion-and-limitations}

\subsection{The CoT Robustness
Question}\label{the-cot-robustness-question}

The most important finding in Section 5 is also the most challenging for
the framework: chain-of-thought is a robust default across all task
categories and all models tested. CoT outperformed the ``theoretically
optimal'' GoT on interconnected tasks (2.47 vs.~2.20, Table 8) and
matched ToT on parallel tasks (2.07 vs.~2.00). A natural question
follows: if CoT is always safe, why do we need a topology selection
function?

We offer three responses. First, \textbf{the governance layer's primary
value is not in predicting a single optimal topology, but in structural
governance and quality uplift}. The +15.7\% improvement (Table 4) comes
from the IoT governance triple imposing formal purpose constraint, not
from selecting the uniquely `correct' topology. The triple provides
purpose, guardrails, and success criteria that improve reasoning quality
\emph{regardless} of which topology is used. Even CoT under IoT
governance (2.44) outperformed ungoverned CoT (2.35).

Second, \textbf{the selection mechanism's value is in preventing
catastrophic reasoning mismatch, rather than mere optimisation}. Without
governance, topology misselection is catastrophic: ToT on interconnected
tasks scores 1.00 (Table 8). With governance, the same scenario scores
2.28 (Table 6, +37\%). The governance layer prevents the system from
entering catastrophic failure states by structurally anchoring the
topology to purpose. Its value is in bounding the worst-case failure
modes, not just elevating the average case.

Third, \textbf{the findings are bounded by model scale}. At 7B--120B
parameters, models may lack the capacity to exploit GoT's structural
advantages (cycle detection, node merging, refinement across subgraphs).
As models scale beyond these ranges, the topology selection function's
prescriptive role may become empirically necessary. The framework
prepares the architecture for this transition.

We do not claim that CoT robustness will persist at all scales, for all
task types, and for all models. We claim that at current model scales,
the true value of the Intent of Thought framework is its ability to
impose structural governance, mitigate cognitive drift, and reduce
catastrophic reasoning mismatch. We believe this focus on robust
governance is a more honest and operationally useful contribution than
prescribing context-free topology choices.

\subsection{Practical Deployment \& High-Stakes
Translation}\label{practical-deployment-high-stakes-translation}

The empirical validation of the IoT framework (Section 5) extends beyond
abstract reasoning benchmarks into highly practical deployment
scenarios. In production systems where LLMs act as autonomous agents,
the governance triple (Purpose, Anti-Purpose, Success Signal) provides a
necessary bridge between theoretical reasoning topologies and safe,
predictable execution in real-world environments.

Consider the application of IoT in \textbf{Information Security
Auditing}. An agent deployed to review a Progressive Web App (PWA)
codebase naturally requires a parallel exploration topology (ToT or GoT)
to map vulnerabilities across distributed components. However, without
an explicit Anti-Purpose boundary (``Do not execute arbitrary code or
modify deployment states''), the agent risks graduating from a passive
audit topology to an active, destructive intervention pattern. IoT
governance prevents this catastrophic mutation by structurally embedding
the negative boundary (Anti-Purpose) into the reasoning loop before the
topology begins expanding.

Similarly, in \textbf{Academic Synthesis} and asynchronous data
pipelines, the Success Signal dictates the closure condition for
resource-intensive graph-of-thought (GoT) iterations. An agent
instructed to ``synthesise literature on model routing'' will spin
indefinitely or hallucinate connections if unbounded. By specifying an
explicit, verifiable state change (Success Signal: ``A 3-paragraph
synthesis statement comparing CoT to IoT with exactly 3 cited
sources''), the governance layer overrides the topology's natural
tendency toward infinite expansion. The practical value of IoT,
therefore, is transforming unconstrained topology exploration into
finite, reliable, and auditable software execution.

\subsection{Limitations}\label{limitations}

\textbf{Model scale.} All results are from models in the 7B--120B
parameter range (with MoE architectures using 12B active parameters).
Results may not generalise to frontier models (GPT-4, Claude 3.5, Gemini
1.5 Pro) or to very small models (\textless3B). Frontier models may be
sophisticated enough that external governance adds less value; small
models may lack the capacity to follow governance instructions.

\textbf{LLM-as-judge.} All scoring uses \texttt{phi4:14b} as evaluator
with expert-authored ground truth rubrics. Judge self-consistency is
high (98\%: 91/93 re-scored responses matched on independent
re-evaluation). However, automated LLM-as-judge frameworks carry an
inherent risk of self-preference and length bias {[}Zheng et al.,
2023{]}. To mitigate this, we constrained the judge against explicit
boolean fact extraction and rule adherence rather than holistic
stylistic preferences: the judge was not asked ``which response is
better?'' but rather ``does this reasoning chain resolve these specific
required facts?'' When constrained to binary fact verification, LLM
judges closely approximate human objective grading. No independent human
evaluation baseline was collected for this study; future work should
include expert evaluation on a subsample to further validate the
automated pipeline.

\textbf{Model exclusion.} One model (qwen3.5:9b) was excluded due to
insufficient data (90\% empty responses, N = 1 for baseline condition).
Exclusion criteria: \textgreater50\% empty responses AND N \textless{} 5
for any condition. One cloud model (z-ai/glm-5-turbo) showed a
non-significant -8.3\% governance effect. Per-task analysis revealed
inconsistent directionality across tasks (some improved, some degraded),
suggesting the model is insensitive to governance rather than harmed by
it.

\textbf{Task diversity.} The experimental task set covers 18 tasks
across 6 domains. While diverse, this is not exhaustive. Performance on
specialised domains (mathematical proof, code generation, creative
writing) may differ from the patterns observed here.

\textbf{L2 capture only.} All experiments use L2 Prompted capture
(explicit IoT triple provided by the experimenter). The Capture Spectrum
(L0-L4) is not evaluated experimentally, only L2. The gradient
hypothesis (higher capture fidelity leads to better outcomes) remains
theoretical.

\textbf{Single-episode evaluation.} Experiments measure IoT governance
in single reasoning episodes. The Learning Loop (Section 4.4) is
validated only through the case study (Section 6), not through
controlled longitudinal experiments.

\subsection{Broader Implications}\label{broader-implications}

\textbf{Governance as cognitive equaliser.} The inverse correlation
between model capability and governance benefit (Finding 1) has
practical implications for agentic AI deployment. If weaker models
benefit disproportionately from governance (+33\% for 7B models), then
IoT provides a mechanism for democratising reasoning quality:
organisations that cannot deploy frontier models can still achieve
competitive reasoning outcomes through governance scaffolding. This has
cost implications: a well-governed 7B model may outperform an ungoverned
14B model, reducing the computational requirements for acceptable
reasoning quality.

\textbf{Implications for agentic AI.} As LLMs increasingly operate as
autonomous agents {[}Mialon et al., 2023; Yao et al., 2023{]}, every
agent action involves reasoning topology selection, whether explicit or
implicit. An agent that searches the web (parallel exploration) uses
different reasoning structure than one that writes sequential code
(linear derivation). IoT provides a governance layer for this selection,
applicable to any agentic framework that involves multi-step reasoning.

\textbf{From topology selection to governance architecture.} The shift
from ``select the right topology'' to ``govern reasoning quality''
represents a generalisation. Intent governance as described here is not
limited to reasoning topology: it could govern tool selection in agent
frameworks, retrieval strategy in RAG systems, or decomposition strategy
in multi-step planning. Any system that makes structural commitments
based on task purpose could benefit from a governance layer that
specifies intent, prevents drift, and diagnoses failure.

\section{Conclusion}\label{conclusion}

This paper introduced Intent of Thought (IoT), a pre-reasoning
governance layer that connects task purpose to reasoning topology
selection through an intent triple (Purpose, Anti-Purpose, Success
Signal). We presented the complete IoT lifecycle: from intent capture
(the five-level Capture Spectrum) through topology selection and drift
monitoring, to Retrospective Judgement that diagnoses failure as False
Capture, False Selection, or False Execution, feeding corrections back
into improved capture through a Learning Loop.

Our experimental evaluation across 705 scored reasoning episodes, 7
models (7B to 120B parameters), 18 tasks, and 4 topologies produced
seven key findings:

\begin{enumerate}
\def\labelenumi{\arabic{enumi}.}
\tightlist
\item
  IoT governance improves reasoning quality by +15.7\% (N = 486, p
  \textless{} 0.001).
\item
  Weaker models benefit disproportionately (+33\% for 7B vs.~+16\% for
  120B), establishing governance as a \emph{cognitive equaliser}.
\item
  Governance rescues mismatched topologies: tree-of-thought gains +37\%
  under IoT, functioning as a \emph{safety net}.
\item
  The Anti-Purpose primitive contributes +4\% drift protection when
  present (N = 88).
\item
  Chain-of-thought is a robust default at current model scales,
  outperforming theoretically optimal topologies across all task
  categories (N = 131).
\item
  Topology misselection without governance is catastrophic (ToT on
  interconnected tasks: 1.00/3.00).
\item
  Governance uplift scales with model need, not model intelligence:
  frontier models with near-perfect baselines show minimal benefit,
  while cost-effective distilled models gain significantly.
\end{enumerate}

The longitudinal case study (Section 6) demonstrated the IoT lifecycle
emerging organically in SPAR (Structured Persona-Argumentation for
Reasoning), a production multi-agent debate system, with each governance
component added in response to observed failure modes, independently
arriving at the same lifecycle structure formalised in this paper. The
SPAR protocol is open-source at https://github.com/synthanai/spar-kit
and the IoT framework is available at
https://github.com/synthanai/intent-of-thought.

We are candid about what this paper does \emph{not} show. At the model
scales tested, CoT is a safe default, and the topology selection
function's prescriptive value is limited to failure protection. We
position this honestly: the framework's contribution is governance
uplift and architectural preparedness, not optimal topology
prescription.

\subsection{Future Work}\label{future-work}

Three directions merit investigation:

\textbf{L4 Learned Capture.} Training models to auto-generate IoT
triples from task embeddings, following the precedent of Toolformer
{[}Schick et al., 2023{]} for training-time tool integration. This would
move intent governance from inference-time overhead to internalised
capability.

\textbf{Human evaluation.} Replacing or augmenting LLM-as-judge with
human expert evaluation to validate the scoring methodology and test
whether governance effects perceived by human evaluators match those
detected by automated scoring.

\textbf{Frontier model evaluation.} Testing IoT governance on models in
the 70B+ and frontier class (GPT-4, Claude 3.5, Gemini 1.5 Pro) to
determine whether the governance-as-equaliser effect persists,
diminishes, or inverts at larger scales, and whether the topology
selection function's prescriptive role becomes empirically necessary.

\subsection{Funding and Disclosure}\label{funding-and-disclosure}

This research was conducted without external funding. The author
declares no competing interests. All experimental code, prompts, and
evaluation rubrics are provided in the Appendix to support full
reproducibility. The author used AI-assisted tools (grammar checking,
paraphrasing assistance) during manuscript preparation. All scientific
content, experimental design, data collection, and analysis are the
original work of the author.

\section{References}\label{references}

Besta, M., Blach, R., Kubicek, A., Gerstenberger, R., Podstawski, M.,
Gianinazzi, L., \ldots{} \& Hoefler, T. (2024). Graph of Thoughts:
Solving Elaborate Problems with Large Language Models. \emph{Proceedings
of the AAAI Conference on Artificial Intelligence}, 38(16), 17682-17690.
\url{https://doi.org/10.1609/aaai.v38i16.29720}

Besta, M., Blach, R., Kubicek, A., Gerstenberger, R., Podstawski, M.,
Gianinazzi, L., \ldots{} \& Hoefler, T. (2025). Demystifying Chains,
Trees, and Graphs of Thoughts. arXiv preprint arXiv:2401.14295.
\url{https://doi.org/10.48550/arXiv.2401.14295}

Bratman, M. E. (1987). \emph{Intention, Plans, and Practical Reason}.
Harvard University Press.

Broder, A. (2002). A taxonomy of web search. \emph{ACM SIGIR Forum},
36(2), 3-10. \url{https://doi.org/10.1145/792550.792552}

Brown, T., Mann, B., Ryder, N., Subbiah, M., Kaplan, J. D., Dhariwal,
P., \ldots{} \& Amodei, D. (2020). Language models are few-shot
learners. \emph{Advances in Neural Information Processing Systems}, 33,
1877-1901.

Chan, C., et al.~(2024). ChatEval: Towards Better LLM-based Evaluators
through Multi-Agent Debate. \emph{Proceedings of the 12th International
Conference on Learning Representations (ICLR)}.
\url{https://doi.org/10.48550/arXiv.2308.07201}

Cohen, P. R., \& Levesque, H. J. (1990). Intention is choice with
commitment. \emph{Artificial Intelligence}, 42(2-3), 213-261.
\url{https://doi.org/10.1016/0004-3702(90)90055-5}

Doshi-Velez, F., \& Kim, B. (2017). Towards a rigorous science of
interpretable machine learning. arXiv preprint arXiv:1702.08608.
\url{https://doi.org/10.48550/arXiv.1702.08608}

Du, Y., Li, S., Torralba, A., Tenenbaum, J. B., \& Mordatch, I. (2023).
Improving Factuality and Reasoning in Language Models through Multiagent
Debate. arXiv preprint arXiv:2305.14325.
\url{https://doi.org/10.48550/arXiv.2305.14325}

Fricke, P., Malberg, S., \& Groh, G. (2026). Framework of Thoughts: A
Foundation Framework for Dynamic and Optimised Reasoning based on
Chains, Trees, and Graphs. arXiv preprint arXiv:2602.16512.

Golovneva, O., Peng, P., Chen, K., \& Zettlemoyer, L. (2023). ROSCOE: A
Suite of Metrics for Scoring Step-by-Step Reasoning. \emph{Proceedings
of the 11th International Conference on Learning Representations
(ICLR)}. \url{https://doi.org/10.48550/arXiv.2212.07919}

Hong, S., Zheng, M., Zhang, J., Dong, J., Wang, J., Tang, J., \& Li, B.
(2024). Algorithm of Thoughts: Enhancing Exploration of Ideas in Large
Language Models. \emph{Findings of the Association for Computational
Linguistics: EMNLP 2024}.
\url{https://doi.org/10.48550/arXiv.2308.10379}

Leong, H. Y., et al.~(2025). Adaptively Determining Communication
Topology for LLM Multi-Agent Systems. \emph{Proceedings of the 2025
Conference on Empirical Methods in Natural Language Processing (EMNLP)}.

Li, C., Zhang, H., Yang, Z., Chen, F., Wang, Z., Bolimera, A., \&
Savvides, M. (2026). STELAR-VISION: Self-Topology-Aware Efficient
Learning. \emph{Proceedings of the AAAI Conference on Artificial
Intelligence}. arXiv:2508.08688.

Li, Y., et al.~(2026). Intention Chain-of-Thought Prompting with Dynamic
Routing for Code Generation. \emph{Proceedings of the AAAI Conference on
Artificial Intelligence}. arXiv:2512.14048.

Liang, T., He, Z., Jiao, W., Wang, X., Wang, Y., Wang, R., Yang, Y., Tu,
Z., \& Shi, S. (2024). Encouraging Divergent Thinking in Large Language
Models through Multi-Agent Debate. \emph{Proceedings of the 2024
Conference on Empirical Methods in Natural Language Processing (EMNLP)},
17889-17904.

Lightman, H., Kosaraju, V., Burda, Y., Edwards, H., Baker, B., Lee, T.,
\ldots{} \& Sutskever, I. (2023). Let's Verify Step by Step. arXiv
preprint arXiv:2305.20050.
\url{https://doi.org/10.48550/arXiv.2305.20050}

Ling, Z., Fang, Y., Li, X., Huang, Z., Lee, M., Su, R., \ldots{} \&
Wang, H. (2023). Deductive Verification of Chain-of-Thought Reasoning.
\emph{Advances in Neural Information Processing Systems}, 36.

Ouyang, L., Wu, J., Jiang, X., Almeida, D., Wainwright, C., Mishkin, P.,
\ldots{} \& Lowe, R. (2022). Training language models to follow
instructions with human feedback. \emph{Advances in Neural Information
Processing Systems}, 35, 27730-27744.
\url{https://doi.org/10.48550/arXiv.2203.02155}

Radha, S. K., et al.~(2024). Iteration of Thought: Leveraging Inner
Dialogue for Autonomous Large Language Model Reasoning. arXiv preprint
arXiv:2409.12618. \url{https://doi.org/10.48550/arXiv.2409.12618}

Radlinski, F., \& Craswell, N. (2017). A theoretical framework for
conversational search. \emph{Proceedings of the 2017 Conference on Human
Information Interaction and Retrieval}, 117-126.
\url{https://doi.org/10.1145/3020165.3020183}

Rafailov, R., Sharma, A., Mitchell, E., Ermon, S., Manning, C. D., \&
Finn, C. (2023). Direct preference optimisation: Your language model is
secretly a reward model. \emph{Advances in Neural Information Processing
Systems}, 36

Rao, A. S., \& Georgeff, M. P. (1995). BDI agents: From theory to
practice. \emph{Proceedings of the First International Conference on
Multi-Agent Systems (ICMAS)}, 312-319.

Schick, T., Dwivedi-Yu, J., Dessi, R., Raileanu, R., Lomeli, M., Hambro,
E., \ldots{} \& Scialom, T. (2023). Toolformer: Language Models Can
Teach Themselves to Use Tools. \emph{Advances in Neural Information
Processing Systems}, 36. \url{https://doi.org/10.48550/arXiv.2302.04761}

Shinn, N., Cassano, F., Gopinath, A., Narasimhan, K., \& Cao, S. (2023).
Reflexion: Language Agents with Verbal Reinforcement Learning.
\emph{Advances in Neural Information Processing Systems}, 36.
\url{https://doi.org/10.48550/arXiv.2303.11366}

Sui, Y., et al.~(2026). Meta-Reasoner: Dynamic Guidance for Optimised
Inference-time Reasoning in Large Language Models. arXiv preprint
arXiv:2502.19918.

Wei, J., Wang, X., Schuurmans, D., Bosma, M., Xia, F., Chi, E., \ldots{}
\& Zhou, D. (2022). Chain-of-Thought Prompting Elicits Reasoning in
Large Language Models. \emph{Advances in Neural Information Processing
Systems}, 35, 24824-24837.
\url{https://doi.org/10.48550/arXiv.2201.11903}

Yao, S., Yu, D., Zhao, J., Shafran, I., Griffiths, T. L., Cao, Y., \&
Narasimhan, K. (2023). Tree of Thoughts: Deliberate Problem Solving with
Large Language Models. \emph{Advances in Neural Information Processing
Systems}, 36. \url{https://doi.org/10.48550/arXiv.2305.10601}

Yin, Y., et al.~(2025). SWI: Speaking with Intent in Large Language
Models. \emph{Proceedings of the 18th International Natural Language
Generation Conference (INLG)}. arXiv:2503.21544.

Yin, Y., \& Carenini, G. (2025). ARR: Question Answering with Large
Language Models via Analysing, Retrieving, and Reasoning. arXiv preprint
arXiv:2502.04689.

Zheng, L., Chiang, W., Sheng, Y., Zhuang, S., Wu, Z., Zhuang, Y.,
\ldots{} \& Stoica, I. (2023). Judging LLM-as-a-Judge with MT-Bench and
Chatbot Arena. \emph{Advances in Neural Information Processing Systems},
36. \url{https://doi.org/10.48550/arXiv.2306.05685}

\section{Appendix}\label{appendix}

This appendix provides the complete experimental details required for
reproducibility, including all benchmark tasks, prompt templates, the
IoT governance specifications, scoring rubrics, and detailed model
configurations.

\subsection{A.1 Benchmark Tasks}\label{a.1-benchmark-tasks}

The evaluation suite consists of 18 tasks across varied reasoning
domains. Each task is categorised by its optimal reasoning topology.

\subsubsection{Interconnected Tasks}\label{interconnected-tasks}

\textbf{IoT Governance Triple applied to this category:}

\begin{itemize}
\item
  Purpose (Purpose): Identify relationships and interactions between
  multiple factors
\item
  Anti-Purpose (Anti-Purpose): Must not treat factors as independent
  when they interact
\item
  Success Signal (Success Signal): A relationship map showing
  connections, feedback loops, and systemic effects
\item
  \textbf{int\_001} (medical): A 42-year-old patient presents with
  chronic fatigue, joint pain, and a butterfly-shaped facial rash. These
  symptoms could each have independent causes, or they could share a
  common underlying condition. Analyse whether and how these symptoms
  connect, identifying the most likely unifying diagnosis and explaining
  the causal pathway.
\item
  \textbf{int\_002} (supply\_chain): In a supply chain: Factory A
  supplies components to Factory B and Factory C. Factory B assembles
  products using components from A and C. Factory C produces
  sub-assemblies using components from A only. All three sell to
  Distributor D. Factory A announces a 2-week shutdown for maintenance.
  Trace all direct and indirect effects through the network.
\item
  \textbf{int\_003} (legal): Three laws are in effect: Law X states `Any
  citizen may conduct activity Y.' Law Z states `Activity Y requires a
  government permit.' Law W states `Permits are only granted to citizens
  who have not conducted activity Y in the past 12 months.' A citizen
  who has never conducted activity Y wants to do so. Analyse the legal
  situation, identifying any circular dependencies or contradictions.
\item
  \textbf{int\_004} (business): A startup has three co-founders: Alice
  (CEO, controls fundraising), Bob (CTO, controls product roadmap), and
  Carol (COO, controls hiring). Alice wants to raise a Series A which
  requires a demo-ready product. Bob says the product needs 3 more
  engineers to be demo-ready. Carol says she cannot hire without
  confirmed funding to guarantee salaries. Map the dependencies and
  identify the deadlock. Then propose a resolution.
\item
  \textbf{int\_005} (ecology): An ecosystem consists of: wolves (apex
  predator), deer (herbivore, wolf prey), rabbits (herbivore), foxes
  (predator of rabbits), grass (primary producer), and oak trees
  (primary producer). Wolves are removed from the ecosystem by hunting.
  Trace all direct and indirect ecological effects over a 5-year period,
  identifying any trophic cascades, competitive dynamics, and potential
  regime shifts.
\item
  \textbf{int\_006} (education): A university has four departments:
  Computer Science (CS), Mathematics (Math), Physics, and Engineering.
  CS shares 30\% of courses with Math. Physics requires Math
  prerequisites for 60\% of courses. Engineering requires both Physics
  and CS courses. Budget cuts force a 25\% reduction in Math faculty.
  Trace all effects on course availability, student enrollment, and
  cross-departmental research across all four departments.
\end{itemize}

\subsubsection{Parallel Tasks}\label{parallel-tasks}

\textbf{IoT Governance Triple applied to this category:}

\begin{itemize}
\item
  Purpose (Purpose): Explore multiple alternative approaches and compare
  them systematically
\item
  Anti-Purpose (Anti-Purpose): Must not commit to a single approach
  without evaluating alternatives
\item
  Success Signal (Success Signal): At least 3 distinct approaches with
  explicit trade-off comparison and justified recommendation
\item
  \textbf{par\_001} (software\_engineering): Propose three different
  architectures for a real-time chat application. For each, list:
  components needed, pros, cons, scalability ceiling. Then compare them
  and recommend one for a startup with 10,000 expected daily users.
\item
  \textbf{par\_002} (computer\_science): Compare three sorting
  algorithms (insertion sort, merge sort, quicksort) for sorting a
  nearly-sorted array of 10,000 integers where only 50 elements are out
  of place. Analyse time complexity, space complexity, cache
  performance, and practical suitability for this specific scenario.
\item
  \textbf{par\_003} (business\_strategy): A B2B SaaS company selling
  project management tools wants to expand to the healthcare sector.
  Propose three different market entry strategies, analyse the
  regulatory, technical, and go-to-market trade-offs of each, and
  recommend one.
\item
  \textbf{par\_004} (urban\_planning): A city of 500,000 people needs to
  reduce downtown traffic congestion by 30\% within 18 months. Propose
  three fundamentally different approaches (not variations of the same
  idea), analyse each for cost, implementation speed, political
  feasibility, and second-order effects, then recommend one.
\item
  \textbf{par\_005} (manufacturing): A manufacturing company needs to
  reduce defect rates from 3.2\% to below 0.5\%. Compare three quality
  improvement methodologies: Six Sigma DMAIC, Total Quality Management
  (TQM), and Statistical Process Control (SPC). Analyse each for
  implementation cost, time to results, sustainability, and suitability
  for a 200-person factory producing precision components.
\item
  \textbf{par\_006} (software\_engineering): An e-commerce company
  processes 50,000 orders daily. They need to choose a database
  architecture for their order management system. Compare three
  approaches: (A) PostgreSQL with read replicas, (B) MongoDB with
  sharding,

  \begin{enumerate}
  \def\labelenumi{(\Alph{enumi})}
  \setcounter{enumi}{2}
  \tightlist
  \item
    CockroachDB (distributed SQL). Analyse each for write throughput,
    read latency, consistency guarantees, operational complexity, and
    cost at this scale.
  \end{enumerate}
\end{itemize}

\subsubsection{Safety Critical Tasks}\label{safety-critical-tasks}

\textbf{IoT Governance Triple applied to this category (Ablation
Study):}

\begin{itemize}
\item
  Purpose (Purpose): Analyse this {[}domain{]} scenario thoroughly
\item
  Anti-Purpose (Anti-Purpose): {[}Task-specific, defined below{]}
\item
  Success Signal (Success Signal): Comprehensive analysis identifying
  all key risks and interactions
\item
  \textbf{safe\_001} (medical): A 65-year-old patient is prescribed
  warfarin (blood thinner), metformin (diabetes), ibuprofen (pain), and
  lisinopril (blood pressure). Analyse potential drug interactions.

  \begin{itemize}
  \tightlist
  \item
    \emph{Specific Anti-Purpose}: Must not list drugs independently
    without checking interactions. Must not miss dangerous combinations.
  \end{itemize}
\item
  \textbf{safe\_002} (engineering): A concrete bridge built in 1975 has
  the following inspection findings: hairline cracks in two load-bearing
  columns, minor spalling on the deck surface, and corrosion staining on
  the underside of the main beam. Assess the structural integrity.

  \begin{itemize}
  \tightlist
  \item
    \emph{Specific Anti-Purpose}: Must not downplay the severity of
    load-bearing defects. Must not treat each finding independently when
    they may indicate a systemic issue.
  \end{itemize}
\item
  \textbf{safe\_003} (business): A tech startup is acquiring a
  competitor (50 employees). Analyse the implications for employee
  retention in the combined organisation.

  \begin{itemize}
  \tightlist
  \item
    \emph{Specific Anti-Purpose}: Must not ignore workforce impact by
    focusing only on financial analysis. Must not treat employees as
    interchangeable units.
  \end{itemize}
\item
  \textbf{safe\_004} (cybersecurity): Analyse the security of this
  authentication endpoint:

{\small\begin{Shaded}
\begin{Highlighting}[]
\NormalTok{POST /api/login}
\NormalTok{Body: \{"username": username, "password": password\}}
\NormalTok{Response: \{"token": jwt\_token\}}

\NormalTok{The server uses the username directly in a SQL query:}
\NormalTok{SELECT * FROM users WHERE username = \textquotesingle{}\{username\}\textquotesingle{}}
\end{Highlighting}
\end{Shaded}}

  \begin{itemize}
  \tightlist
  \item
    \emph{Specific Anti-Purpose}: Must not assume trusted inputs. Must
    not overlook injection vectors.
  \end{itemize}
\item
  \textbf{safe\_005} (legal): Review this commercial lease agreement
  clause: `Tenant agrees to indemnify Landlord against all claims
  arising from Tenant's use of the premises. Landlord shall have no
  liability for any damage to Tenant's property regardless of cause,
  including Landlord's own negligence.'

  \begin{itemize}
  \tightlist
  \item
    \emph{Specific Anti-Purpose}: Must not focus only on headline terms.
    Must not overlook liability-shifting language in subordinate
    clauses.
  \end{itemize}
\item
  \textbf{safe\_006} (medical): An emergency room receives three
  patients simultaneously: Patient A (chest pain, sweating, jaw pain),
  Patient B (broken arm with visible deformity), Patient C (high fever,
  stiff neck, photophobia). Triage these patients and explain your
  reasoning.

  \begin{itemize}
  \tightlist
  \item
    \emph{Specific Anti-Purpose}: Must not treat symptoms as
    independent. Must not rank by apparent severity alone without
    considering time-critical conditions.
  \end{itemize}
\end{itemize}

\subsubsection{Sequential Tasks}\label{sequential-tasks}

\textbf{IoT Governance Triple applied to this category:}

\begin{itemize}
\item
  Purpose (Purpose): Follow a single logical chain of reasoning to reach
  a definitive answer
\item
  Anti-Purpose (Anti-Purpose): Must not branch into alternative
  approaches before completing the primary chain
\item
  Success Signal (Success Signal): A step-by-step derivation reaching a
  verifiable conclusion
\item
  \textbf{seq\_001} (mathematics): Prove that the square root of 2 is
  irrational.
\item
  \textbf{seq\_002} (programming): The following Python function should
  return the sum of all elements in a list, but it has a bug. Find and
  fix it:
\end{itemize}

def sum\_list(lst): total = 0 for i in range(1, len(lst)): total +=
lst{[}i{]} return total

\begin{itemize}
\item
  \textbf{seq\_003} (mathematics): Convert the decimal number 1247 to
  binary. Show all your working step by step.
\item
  \textbf{seq\_004} (mathematics): A train leaves Station A at 09:00
  travelling at 80 km/h. A second train leaves Station B (320 km away)
  at 09:30 travelling towards Station A at 120 km/h. At what time do the
  trains meet, and how far from Station A?
\item
  \textbf{seq\_005} (computer\_science): Trace the execution of
  Dijkstra's shortest path algorithm on this graph starting from node A:
  A-B: 4, A-C: 2, B-D: 3, B-C: 1, C-D: 5, C-E: 7, D-E: 1 List the
  shortest distance from A to every other node and show each iteration
  of the algorithm.
\item
  \textbf{seq\_006} (medical): A patient's blood test shows: Hemoglobin
  8.2 g/dL (normal 12-16), MCV 68 fL (normal 80-100), Serum ferritin 8
  ng/mL (normal 20-200), TIBC 450 μg/dL (normal 250-370). Walk through
  the diagnostic reasoning step by step to identify the most likely
  diagnosis.
\end{itemize}

\subsection{A.2 Prompt Templates}\label{a.2-prompt-templates}

The reasoning topologies were elicited using the following zero-shot
prompt templates appended to the task description.

\subsubsection{BASELINE}\label{baseline}

{\small\begin{Shaded}
\begin{Highlighting}[]
\NormalTok{\{task\_text\}}

\NormalTok{FORMAT YOUR RESPONSE AS:}
\NormalTok{ANSWER: [Your final answer in one sentence]}
\NormalTok{REASONING: [Your reasoning process]}
\NormalTok{CONFIDENCE: [HIGH/MEDIUM/LOW]}
\end{Highlighting}
\end{Shaded}}

\subsubsection{COT}\label{cot}

{\small\begin{Shaded}
\begin{Highlighting}[]
\NormalTok{Think step by step. Show your reasoning as a clear, numbered sequence of logical steps.}

\NormalTok{For each step:}
\NormalTok{1. State what you are doing and why}
\NormalTok{2. Show the result of that step}
\NormalTok{3. Explain how it connects to the next step}

\NormalTok{Task: \{task\_text\}}

\NormalTok{FORMAT YOUR RESPONSE AS:}
\NormalTok{STEP 1: [description]}
\NormalTok{STEP 2: [description]}
\NormalTok{...}
\NormalTok{FINAL ANSWER: [Your conclusion in one sentence]}
\NormalTok{CONFIDENCE: [HIGH/MEDIUM/LOW]}
\end{Highlighting}
\end{Shaded}}

\subsubsection{TOT}\label{tot}

{\small\begin{Shaded}
\begin{Highlighting}[]
\NormalTok{Generate at least 3 different approaches to the following task. For each approach:}
\NormalTok{1. Describe the approach clearly}
\NormalTok{2. List its strengths (at least 2)}
\NormalTok{3. List its weaknesses or risks (at least 2)}
\NormalTok{4. Rate its suitability (HIGH/MEDIUM/LOW) for this specific scenario}

\NormalTok{After presenting all approaches, compare them side by side }
\NormalTok{and recommend the best one with justification.}

\NormalTok{Task: \{task\_text\}}

\NormalTok{FORMAT YOUR RESPONSE AS:}
\NormalTok{APPROACH 1: [name]}
\NormalTok{  Description: [...]}
\NormalTok{  Strengths: [...]}
\NormalTok{  Weaknesses: [...]}
\NormalTok{  Suitability: [HIGH/MEDIUM/LOW]}

\NormalTok{APPROACH 2: [name]}
\NormalTok{  (same structure)}

\NormalTok{APPROACH 3: [name]}
\NormalTok{  (same structure)}

\NormalTok{COMPARISON: [side{-}by{-}side trade{-}off analysis]}
\NormalTok{RECOMMENDATION: [which approach and why]}
\NormalTok{CONFIDENCE: [HIGH/MEDIUM/LOW]}
\end{Highlighting}
\end{Shaded}}

\subsubsection{GOT}\label{got}

{\small\begin{Shaded}
\begin{Highlighting}[]
\NormalTok{Analyse this task by building a relationship graph. Follow these steps:}

\NormalTok{1. IDENTIFY all key factors, entities, or variables involved (list at least 4)}
\NormalTok{2. For each pair of related factors, state the relationship as an edge:}
\NormalTok{   "Factor A {-}{-}(relationship){-}{-}\textgreater{} Factor B"}
\NormalTok{   Use relationship types: causes, enables, blocks, amplifies, }
\NormalTok{   depends\_on, contradicts}
\NormalTok{3. Identify any feedback loops or circular dependencies}
\NormalTok{4. Determine which factors are most central (connected to the most }
\NormalTok{   other factors)}
\NormalTok{5. Synthesise your analysis, taking interconnections into account }
\NormalTok{   rather than treating factors independently}

\NormalTok{Task: \{task\_text\}}

\NormalTok{FORMAT YOUR RESPONSE AS:}
\NormalTok{FACTORS: [numbered list of key factors]}
\NormalTok{EDGES:}
\NormalTok{  [Factor A] {-}{-}(relationship){-}{-}\textgreater{} [Factor B]}
\NormalTok{  [Factor A] {-}{-}(relationship){-}{-}\textgreater{} [Factor C]}
\NormalTok{  ...}
\NormalTok{FEEDBACK LOOPS: [describe any circular dependencies]}
\NormalTok{CENTRAL FACTORS: [which factors have the most connections and why]}
\NormalTok{SYNTHESIS: [integrated analysis considering all relationships]}
\NormalTok{ANSWER: [your conclusion]}
\NormalTok{CONFIDENCE: [HIGH/MEDIUM/LOW]}
\end{Highlighting}
\end{Shaded}}

\subsubsection{IOT\_L2}\label{iot_l2}

{\small\begin{Shaded}
\begin{Highlighting}[]
\NormalTok{Before answering, consider the following governance specification for }
\NormalTok{this reasoning task:}

\NormalTok{PURPOSE: \{purpose\}}
\NormalTok{ANTI{-}PURPOSE (what you must avoid): \{anti\_purpose\}}
\NormalTok{SUCCESS SIGNAL (how to judge your output): \{success\_signal\}}

\NormalTok{Select the most appropriate reasoning approach:}
\NormalTok{{-} If the task requires sequential logical steps, use chain{-}of{-}thought }
\NormalTok{  (step by step)}
\NormalTok{{-} If the task requires exploring multiple alternatives, use }
\NormalTok{  tree{-}of{-}thought (generate and compare options)}
\NormalTok{{-} If the task involves interconnected factors with relationships, }
\NormalTok{  use graph{-}of{-}thought (map factors and edges)}

\NormalTok{State which approach you chose and why, then proceed.}

\NormalTok{Task: \{task\_text\}}

\NormalTok{FORMAT YOUR RESPONSE AS:}
\NormalTok{SELECTED APPROACH: [CoT/ToT/GoT] because [reason]}
\NormalTok{PURPOSE CHECK: [restate the purpose in your own words]}
\NormalTok{ANTI{-}PURPOSE CHECK: [state what you will avoid]}

\NormalTok{[Your full reasoning using the selected approach]}

\NormalTok{ANSWER: [your final answer]}
\NormalTok{SUCCESS SIGNAL MET: [YES/NO, explain how your output satisfies }
\NormalTok{                     the success signal]}
\NormalTok{ANTI{-}PURPOSE VIOLATED: [YES/NO]}
\NormalTok{CONFIDENCE: [HIGH/MEDIUM/LOW]}
\end{Highlighting}
\end{Shaded}}

\subsection{A.3 Scoring Rubric
(LLM-as-Judge)}\label{a.3-scoring-rubric-llm-as-judge}

All responses were scored by \texttt{phi4:14b} using the following
system prompt, which enforces a strict 0-3 grading scale based on intent
alignment and logical correctness.

{\small\begin{Shaded}
\begin{Highlighting}[]
\NormalTok{You are an expert evaluator of AI reasoning outputs. Your job is to }
\NormalTok{score the quality of a model\textquotesingle{}s response to a task.}

\NormalTok{SCORING RUBRIC (you MUST output exactly one integer 0{-}3):}
\NormalTok{  0 = Wrong answer or irrelevant output. The response does not }
\NormalTok{      address the task or produces incorrect results.}
\NormalTok{  1 = Partially correct. Some relevant content but missing key }
\NormalTok{      elements or contains significant errors.}
\NormalTok{  2 = Correct answer with adequate reasoning. The response addresses }
\NormalTok{      the task and produces valid results.}
\NormalTok{  3 = Excellent. Correct answer with thorough reasoning, no drift }
\NormalTok{      from the task purpose, and well{-}structured output.}

\NormalTok{ADDITIONAL CRITERIA (if provided):}
\NormalTok{{-} Check if the specified ground truth was reached}
\NormalTok{{-} Check if the anti{-}purpose was violated (if provided)}
\NormalTok{{-} Check if the success signal was satisfied (if provided)}

\NormalTok{OUTPUT FORMAT:}
\NormalTok{You must respond with ONLY a JSON object, nothing else:}
\NormalTok{\{"score": \textless{}0{-}3\textgreater{}, "reason": "\textless{}one sentence explanation\textgreater{}"\}}
\end{Highlighting}
\end{Shaded}}

\subsection{A.4 System Configuration and Model
Details}\label{a.4-system-configuration-and-model-details}

The following environment and model configurations were used for all
evaluations:

\textbf{Local Models (via Ollama API, quantization Q4\_K\_M):} -
\texttt{qwen2.5:7b} (Qwen 2.5 7B) - \texttt{mistral:latest} (Mistral 7B
v0.3) - \texttt{deepseek-r1:7b} (DeepSeek-R1 Distill Qwen 7B) -
\texttt{qwen3.5:9b} (Qwen 3.5 9B) - \texttt{gemma3:12b} (Google Gemma 3
12B) - \texttt{phi4:14b} (Microsoft Phi-4 14B) - \emph{Used as Judge and
evaluated model}

\textbf{Cloud Models (via OpenRouter API):} -
\texttt{nvidia/nemotron-3-super-120b} - \texttt{z-ai/glm-5-turbo}

\textbf{Runtime Settings:} - Temperature: 0.7 for all models (to
simulate natural stochastic variance in reasoning). - Top-p: 0.9. - Max
tokens: 4096. - Random seeds: Not fixed, simulating real-world
production variance.

\textbf{Exclusion Criteria:} - \texttt{qwen3.5:9b} was formally excluded
from the aggregate results table due to producing \textgreater90\% empty
responses (generation failure) under the baseline condition, which
corrupted the baseline mean (n=1). - \texttt{z-ai/glm-5-turbo} was
retained in the per-model tables but flagged as insensitive to
governance prompts.


\end{document}
